\let\R\relax
\newcommand{\R}{\mathbb{R}}
\let\C\relax
\newcommand{\C}{\mathbb{C}}
\let\N\relax
\newcommand{\N}{\mathbb{N}}
\let\Z\relax
\newcommand{\Z}{\mathbb{Z}}
\let\WF\relax
\newcommand{\WF}{\operatorname{WF}}
\let\singsupp\relax
\newcommand{\singsupp}{\operatorname{singsupp}}

\chapter{Lars H\"{o}rmander (1988)}

\begin{quote}
\it ``H\"ormander's profound and wide-ranging contributions to the modern
theory of linear partial differential equations have transformed the
subject. His monumental four-volume treatise \textit{The Analysis of Linear
Partial Differential Operators} is the definitive account of the field,
synthesizing pseudodifferential operators, Fourier integral operators,
and microlocal analysis into a unified theory of unprecedented power and
elegance.'' --- Wolf Prize Citation
\end{quote}

Lars H\"ormander (1931--2012) was the preeminent analyst of partial
differential equations in the second half of the twentieth century. The
1988 Wolf Prize in Mathematics (shared with Friedrich Hirzebruch)
recognized his transformation of the theory of linear PDEs through the
development of microlocal analysis, pseudodifferential operators,
Fourier integral operators, and the systematic study of local solvability
and hypoellipticity. His four-volume magnum opus, \textit{The Analysis of
Linear Partial Differential Operators} (1983--1985), remains the
definitive reference and the crowning achievement of the field.

H\"ormander's work is characterized by extraordinary technical mastery,
an insistence on the most general formulations, and a drive to reduce
the local theory of linear PDEs to algebraic and geometric conditions
that can be verified by computation. Before H\"ormander, the theory of
linear PDEs was a collection of specialized results for different types
of equations (elliptic, hyperbolic, parabolic). After H\"ormander, it
became a unified subject based on the microlocal analysis of singularities,
where the symbol of an operator and the geometry of its characteristic
set determine the behavior of solutions.

\section{Main Contribution}

\begin{itemize}
    \item \textbf{Pseudodifferential Operators (1965--1971):} H\"ormander
    developed the complete theory of pseudodifferential operators,
    generalizing differential operators to operators with symbols that
    are smooth functions on the cotangent bundle. He classified operators
    by their symbol classes \(S^m_{\rho,\delta}\) and established their
    algebraic and analytic properties, including composition rules,
    parametrices, and \(L^2\) boundedness. This theory provides a unified
    framework for constructing approximate inverses (parametrices) of
    elliptic and hypoelliptic operators.
    \item \textbf{Fourier Integral Operators (1971):} H\"ormander
    generalized pseudodifferential operators to Fourier integral
    operators, which are oscillatory integrals whose phase functions are
    not necessarily linear. These operators are the natural tools for
    solving hyperbolic equations and describing the propagation of
    singularities. The Fourier integral operator calculus is the
    rigorous mathematical foundation for geometrical optics.
    \item \textbf{Wave Front Sets and Microlocal Analysis (1970--1971):}
    H\"ormander introduced the concept of the wave front set
    \(\WF(u)\) of a distribution \(u\), which refines the notion of
    singular support by recording not only \textit{where} a distribution
    is singular but also \textit{which frequencies} cause the
    singularity. This is the central concept of microlocal analysis, the
    study of distributions and PDEs in the cotangent bundle.
    \item \textbf{Propagation of Singularities (1970--1971):}
    H\"ormander proved that for a real principal-type operator \(P\),
    the wave front set of a solution \(u\) to \(Pu = f\) propagates along
    the bicharacteristics of the principal symbol. This theorem
    generalizes the classical fact that singularities of solutions of
    the wave equation travel along light rays.
    \item \textbf{H\"ormander Condition for Local Solvability (1960s):}
    H\"ormander gave a necessary and sufficient condition for a linear
    partial differential operator to be locally solvable: the principal
    symbol must not change sign when crossing a characteristic surface
    of double multiplicity. This resolved a long-standing problem
    originating from Hans Lewy's 1957 example of a non-solvable
    equation.
    \item \textbf{Hypoellipticity (1960s):} H\"ormander characterized
    hypoelliptic operators --- operators \(P\) such that \(Pu \in
    C^\infty\) implies \(u \in C^\infty\) --- in terms of a condition on
    the iterated commutators of the vector fields defining \(P\). This
    is the fundamental theorem on sums of squares of vector fields.
    \item \textbf{H\"ormander Spaces (1960s):} H\"ormander introduced
    the function spaces \(B_{p,k}\) and \(H_{(s)}\) (Sobolev spaces with
    variable order of regularity), providing the functional-analytic
    framework for the sharp regularity theory of linear PDEs. These
    spaces are now called H\"ormander spaces or H\"ormander--Ljubich
    spaces.
    \item \textbf{Four-Volume Treatise (1983--1985):} \textit{The
    Analysis of Linear Partial Differential Operators} (Volumes 1--4)
    is the definitive treatment of the subject, synthesizing every major
    development from distribution theory through Fourier integral
    operators into a coherent, self-contained exposition.
\end{itemize}

\section{Origin of the Problem}

\subsection{The Solvability Crisis}

In 1957, Hans Lewy shocked the mathematical world by exhibiting a linear
partial differential equation with smooth coefficients that has no
solution --- even in the sense of distributions --- for any choice of
smooth right-hand side. The equation was
\[
L u = \frac{\partial u}{\partial x_1} + i \frac{\partial u}{\partial x_2}
- 2i(x_1 + i x_2) \frac{\partial u}{\partial x_3} = f.
\]
Lewy showed that for most smooth \(f\), there is no distribution \(u\)
satisfying \(Lu = f\) in any open set.

This discovery shattered the long-held belief that every linear PDE with
smooth coefficients has solutions, at least locally. The problem of
characterizing which operators are \textit{locally solvable} --- for
which every smooth right-hand side admits a local distribution
solution --- became one of the central problems of linear PDE theory.

\subsection{The Elliptic Paradigm and Its Limits}

Before H\"ormander's work, the theory of linear PDEs was organized
around the classification into elliptic, hyperbolic, and parabolic
types. Elliptic operators (like the Laplacian) have the best properties:
solutions are smooth wherever the coefficients and right-hand side are
smooth, and parametrices (approximate inverses) can be constructed using
the theory of singular integrals. Hyperbolic operators (the wave
operator) propagate singularities along characteristics. Parabolic
operators (the heat operator) instantaneously smooth solutions.

These classifications, however, were based on the principal part of the
operator and its algebraic properties. They did not capture finer
phenomena: there exist operators that are neither elliptic nor
hypoelliptic in the classical sense, yet have interesting regularity
properties. The need for a more refined, microlocal approach --- one
that examines operators and their solutions in phase space (position
and frequency simultaneously) --- became increasingly clear.

\subsection{The Problem of Singularities}

A fundamental question in linear PDE theory is: given an operator \(P\)
and a solution \(u\) to \(Pu = f\), how are the singularities of \(u\)
related to those of \(f\) and the geometry of the operator? For
elliptic operators, \(\singsupp(u) = \singsupp(f)\) (elliptic regularity).
For the wave equation, singularities propagate along light rays. But
what about general operators?

This question requires a precise language for describing singularities.
The classical notion of singular support --- the complement of the set
where a distribution is smooth --- is too coarse. It tells us
\textit{where} the singularities are, but not \textit{what direction}
they come from. A theory of \textit{directional} or
\textit{microlocal} singularities was needed.

\subsection{The Challenge of Hypoellipticity}

An operator \(P\) is called \textit{hypoelliptic} if \(Pu \in C^\infty\)
implies \(u \in C^\infty\). Elliptic operators are hypoelliptic, but the
converse is false. The classic example is the heat operator
\(\partial_t - \partial_x^2\), which is not elliptic but is hypoelliptic
(it smooths initial data instantaneously).

The problem of characterizing hypoellipticity in geometric terms was
open. The Kolmogorov--Fokker--Planck equation,
\[
\partial_t u - \partial_x^2 u - x \partial_y u = 0,
\]
is hypoelliptic even though its symbol vanishes on a large set.
Understanding why requires analyzing the commutation relations of the
vector fields involved.

\section{How It Was Solved}

\subsection{The H\"ormander Condition for Local Solvability}

Let \(P(x,D)\) be a linear partial differential operator of order \(m\)
with smooth coefficients and principal symbol \(p_m(x,\xi)\). The
operator is of \textit{principal type} if \(dp_m \neq 0\) whenever
\(p_m = 0\) and \(\xi \neq 0\).

\begin{definition}[Local Solvability]
An operator \(P\) is \textbf{locally solvable} at a point \(x_0\) if
there exists a neighborhood \(\Omega\) of \(x_0\) such that for every
\(f \in C_0^\infty(\Omega)\), there exists a distribution
\(u \in \mathcal{D}'(\Omega)\) satisfying \(Pu = f\) in \(\Omega\).
\end{definition}

\begin{theorem}[H\"ormander's Local Solvability Condition, 1960]
Let \(P\) be a linear partial differential operator of principal type
with smooth coefficients. Then \(P\) is locally solvable at \(x_0\) if
and only if the following condition holds: for every characteristic
point \((x_0, \xi^0)\) with \(p_m(x_0, \xi^0) = 0\), \(\xi^0 \neq 0\),
and for every real-valued function \(\phi\) such that \(\phi(x_0) = 0\),
\(d\phi(x_0) = \xi^0\), the function
\[
\tau \mapsto \overline{p_m(x, \xi + i\tau\,d\phi(x))}
\]
does not change sign as \(\tau\) passes through zero for \(x\) near
\(x_0\) and \(\xi\) near \(\xi^0\).
\end{theorem}

In simpler language, the condition requires that when the operator is
written in a coordinate system where one coordinate is the
characteristic direction, the principal symbol must not change sign
when the frequency variable corresponding to that direction crosses the
characteristic surface. This condition is often expressed as: \(p_m\)
must satisfy the \textit{sign condition} that for every real zero
\((x,\xi)\) of \(p_m\), the imaginary part of \(p_m\) (when the
frequency is given a small imaginary increment) does not vanish to
infinite order.

This theorem resolved the problem raised by Lewy's example: Lewy's
operator fails the sign condition, which is why it has no local
solutions. For operators with real principal symbols (such as elliptic,
hyperbolic, and parabolic operators), the condition is automatically
satisfied, confirming that these operators are locally solvable.

\subsection{Wave Front Sets and Microlocal Regularity}

H\"ormander's microlocal analysis begins with a precise notion of
directional singularities. The key idea is to examine a distribution
not just in physical space but in the cotangent bundle \(T^*X\) of a
manifold \(X\), where the fiber directions represent frequencies.

\begin{definition}[Wave Front Set]
Let \(u \in \mathcal{D}'(X)\) be a distribution on an open set
\(X \subset \R^n\). A point \((x_0, \xi_0) \in T^*X \setminus 0\) is
called a \textbf{regular direction} for \(u\) if there exists a
cutoff function \(\phi \in C_0^\infty(X)\) with \(\phi(x_0) \neq 0\)
and a conic neighborhood \(\Gamma\) of \(\xi_0\) such that
\[
|\widehat{\phi u}(\xi)| \leq C_N (1+|\xi|)^{-N}
\]
for all \(N\) and all \(\xi \in \Gamma\). The \textbf{wave front set}
\(\WF(u) \subset T^*X \setminus 0\) is the complement of the set of
regular directions.
\end{definition}

In other words, \((x_0, \xi_0) \notin \WF(u)\) means that near \(x_0\),
the distribution \(u\) is smooth in the direction \(\xi_0\). The wave
front set refines the singular support:
\[
\pi(\WF(u)) = \singsupp(u),
\]
where \(\pi: T^*X \to X\) is the projection.

\begin{example}[Wave Front Set of the Delta Function]
The Dirac delta \(\delta_0\) at the origin in \(\R^n\) has
\[
\WF(\delta_0) = \{(0, \xi) : \xi \in \R^n \setminus 0\},
\]
since \(\widehat{\phi\delta_0}(\xi) = \phi(0)\) does not decay in any
direction.
\end{example}

\begin{example}[Wave Front Set of a Step Function]
The Heaviside function \(H(x_1)\) on \(\R^n\) has
\[
\WF(H) = \{(x, \xi) : x_1 = 0, \xi_1 \neq 0, \xi_2 = \cdots = \xi_n = 0\},
\]
reflecting that the singularity has a specific direction (normal to
the discontinuity).
\end{example}

The wave front set is the natural language for describing operations on
distributions. The product \(uv\) of two distributions is well-defined
if \(\WF(u) + \WF(v)\) does not contain any zero fibers (H\"ormander's
criterion for multiplication of distributions).

\subsection{Propagation of Singularities}

The most celebrated theorem in microlocal analysis describes how
singularities move under the action of a differential operator.

Let \(P\) be a pseudodifferential operator of order \(m\) with principal
symbol \(p_m(x,\xi)\) which is real-valued and homogeneous of degree
\(m\) in \(\xi\). The \textit{characteristic set} is
\[
\Sigma = \{(x,\xi) \in T^*X \setminus 0 : p_m(x,\xi) = 0\}.
\]
The \textit{Hamilton vector field} of \(p_m\) is
\[
H_{p_m} = \sum_{j=1}^n \left(\frac{\partial p_m}{\partial \xi_j}
\frac{\partial}{\partial x_j} - \frac{\partial p_m}{\partial x_j}
\frac{\partial}{\partial \xi_j}\right).
\]
The integral curves of \(H_{p_m}\) on \(\Sigma\) are called
\textit{bicharacteristics}.

\begin{theorem}[Propagation of Singularities, H\"ormander 1970--1971]
Let \(P\) be a pseudodifferential operator of real principal type on a
manifold \(X\). Suppose \(u \in \mathcal{D}'(X)\) satisfies
\(Pu = f\) with \(f\) smooth. Then the wave front set of \(u\) is
contained in the characteristic set \(\Sigma\) and is invariant under
the Hamilton flow:
\[
\WF(u) \setminus \WF(f) \subset \Sigma,
\]
and for every bicharacteristic curve \(\gamma\) of \(p_m\),
\[
\gamma \cap (\WF(u) \setminus \WF(f)) \neq \varnothing
\;\Longrightarrow\; \gamma \subset \WF(u) \setminus \WF(f).
\]
In particular, if \(u\) is smooth at one point of a bicharacteristic,
it is smooth along the entire bicharacteristic, modulo the smoothing
effect of the right-hand side.
\end{theorem}

This theorem generalizes the classical fact that for the wave equation
\(\square u = 0\), singularities propagate along light rays. For the
wave operator \(\square\), the principal symbol is
\(p_2(x,\xi) = \xi_0^2 - |\xi'|^2\), the characteristic set is the
light cone, and the bicharacteristics are precisely the light rays.

\begin{proof}[Idea of Proof]
The proof uses the calculus of pseudodifferential operators and
microlocal parametrices. One constructs a microlocal inverse (a
parametrix) for \(P\) near a non-characteristic point, then analyzes
its action on the wave front set. Near a characteristic point, one
factors \(P\) into first-order factors and uses energy estimates along
bicharacteristics. The key technical ingredient is the construction of
a pseudodifferential operator \(A\) whose symbol is supported near a
bicharacteristic segment, such that \([P, A]\) is of lower order,
allowing one to localize the propagation estimate.
\end{proof}

\subsection{Pseudodifferential Operators}

A pseudodifferential operator is a generalization of a differential
operator that allows the symbol to be a smooth function on the
cotangent bundle rather than a polynomial in \(\xi\).

\begin{definition}[Symbol Class \(S^m_{\rho,\delta}\)]
For \(m \in \R\) and \(0 \leq \delta < \rho \leq 1\), the symbol class
\(S^m_{\rho,\delta}(\R^n \times \R^n)\) consists of smooth functions
\(a(x,\xi)\) such that for all multi-indices \(\alpha, \beta\),
\[
|\partial_x^\beta \partial_\xi^\alpha a(x,\xi)|
\leq C_{\alpha,\beta} (1+|\xi|)^{m - \rho|\alpha| + \delta|\beta|}.
\]
The class \(S^m_{1,0}\) is the most commonly used, often denoted simply
by \(S^m\).
\end{definition}

A pseudodifferential operator \(a(x,D)\) with symbol \(a(x,\xi)\) acts
on a Schwarz function \(u\) by the oscillatory integral
\[
a(x,D) u(x) = (2\pi)^{-n} \int_{\R^n} e^{i x \cdot \xi}
a(x,\xi) \hat u(\xi) \,d\xi.
\]

\begin{theorem}[Composition of Pseudodifferential Operators, H\"ormander 1965]
If \(a \in S^{m_1}_{\rho,\delta}\) and \(b \in S^{m_2}_{\rho,\delta}\),
then the composition \(a(x,D) \circ b(x,D)\) is a pseudodifferential
operator with symbol \(c \in S^{m_1+m_2}_{\rho,\delta}\) given by the
asymptotic expansion
\[
c(x,\xi) \sim \sum_{\alpha} \frac{1}{\alpha!}
\partial_\xi^\alpha a(x,\xi) \, D_x^\alpha b(x,\xi),
\]
where \(D_x = -i\partial_x\).
\end{theorem}

This composition formula is fundamental. Its leading term is simply the
product of symbols, and the lower-order terms involve derivatives of the
symbols. The proof uses the stationary phase method and the theory of
oscillatory integrals.

The existence of parametrices for elliptic operators follows directly:
if \(p_m(x,\xi)\) is the principal symbol of an elliptic operator \(P\),
then \(a(x,\xi) = 1/p_m(x,\xi)\) is a symbol of class \(S^{-m}\), and
\(a(x,D)\) is a parametrix for \(P\), meaning \(P \circ a(x,D) = I + R\)
where \(R\) is a smoothing operator.

\subsection{Fourier Integral Operators}

Fourier integral operators generalize pseudodifferential operators by
allowing more general phase functions. While a pseudodifferential
operator has phase function \(\phi(x,\xi) = x \cdot \xi\), a Fourier
integral operator allows any phase function that satisfies certain
non-degeneracy conditions.

\begin{definition}[Fourier Integral Operator]
A Fourier integral operator \(A\) is defined by
\[
A u(x) = (2\pi)^{-n} \iint_{\R^n \times \R^N}
e^{i\phi(x,\theta)} a(x,\theta) \hat u(\xi) \,d\xi\,d\theta,
\]
or more generally by
\[
A u(x) = \int_{\R^N} e^{i\phi(x,\theta)} a(x,\theta) \hat u(\xi) \,d\theta,
\]
where \(\phi\) is a non-degenerate phase function and
\(a \in S^m_{\rho,\delta}\) is a symbol. The operator \(A\) is
associated with a canonical relation \(\Lambda \subset T^*X \times T^*Y\)
given by
\[
\Lambda = \{(x, d_x\phi, y, -d_y\phi) : d_\theta\phi = 0\}.
\]
\end{definition}

Fourier integral operators are the natural tools for solving hyperbolic
equations. The solution operator for the wave equation,
\[
u(t, x) = \int_{\R^n} e^{i(x \cdot \xi \pm t|\xi|)}
\hat u_0(\xi) \,d\xi + \text{(lower order terms)},
\]
is a Fourier integral operator whose canonical relation is the
bicharacteristic flow of the wave operator.

\begin{theorem}[Composition of Fourier Integral Operators, H\"ormander 1971]
Fourier integral operators are closed under composition when the
canonical relations are composable. If \(A\) and \(B\) are Fourier
integral operators associated with canonical relations \(\Lambda_A\)
and \(\Lambda_B\), then \(A \circ B\) (when defined) is a Fourier
integral operator associated with the composition
\[
\Lambda_{A \circ B} = \Lambda_A \circ \Lambda_B,
\]
provided the composition is clean (a transversality condition on the
intersection of the two relations).
\end{theorem}

The theory of Fourier integral operators provides a rigorous foundation
for geometrical optics and the WKB method. The leading term of the
oscillatory integral solution corresponds to the transport equation
along rays, while the lower-order terms account for diffraction and
other wave phenomena.

\subsection{Hypoellipticity and Sums of Squares}

An operator \(P\) is hypoelliptic if \(Pu \in C^\infty\) forces
\(u \in C^\infty\). Equivalently, the singular support of \(u\) equals
the singular support of \(Pu\). H\"ormander gave a geometric condition
for hypoellipticity in terms of vector fields.

\begin{definition}[H\"ormander's Condition on Vector Fields]
Let \(X_1, \ldots, X_k\) be smooth vector fields on an open set
\(\Omega \subset \R^n\). These vector fields satisfy
\textbf{H\"ormander's condition} of step \(r\) if the Lie brackets
\[
\{X_{i_1}, [X_{i_2}, [\cdots, X_{i_j}]\cdots] : 1 \leq j \leq r\}
\]
span the tangent space at every point.
\end{definition}

\begin{theorem}[H\"ormander's Hypoellipticity Theorem, 1967]
Let \(X_0, X_1, \ldots, X_k\) be smooth vector fields satisfying
H\"ormander's condition on an open set \(\Omega \subset \R^n\). Then
the operator
\[
P = X_0 + \sum_{j=1}^k X_j^2
\]
is hypoelliptic. More precisely, if \(u\) is a distribution on
\(\Omega\) and \(Pu \in H^s_{\text{loc}}(\Omega)\), then
\(u \in H^{s+\varepsilon}_{\text{loc}}(\Omega)\) for some
\(\varepsilon > 0\) depending on the step of the bracket condition.
\end{theorem}

This theorem explains the hypoellipticity of the heat operator
\(\partial_t - \partial_x^2\): the vector fields \(X_0 = \partial_t\) and
\(X_1 = \partial_x\) bracket to zero (they commute), so the step is 1.
For the Kolmogorov equation \(\partial_t - \partial_x^2 - x\partial_y\),
the vector fields are
\[
X_0 = \partial_t, \quad X_1 = \partial_x, \quad X_2 = x\partial_y,
\]
and their commutator is \([X_1, X_2] = \partial_y\), which together with
\(X_1, X_2\) spans three directions, giving hypoellipticity with
step 2.

More generally, if the operator is of the form
\[
P = \sum_{j=1}^k X_j^2 + X_0 + c,
\]
and the Lie algebra generated by the vector fields has full rank at
every point, then \(P\) is hypoelliptic with a loss of derivatives
determined by the step of the Lie algebra. This result is sometimes
called the \textit{sum of squares theorem}.

\subsection{H\"ormander Spaces}

H\"ormander introduced function spaces that allow the order of
regularity to vary from point to point and from direction to direction.
These are Sobolev spaces with variable order.

\begin{definition}[H\"ormander Space \(H_{(s)}\)]
Let \(s : T^*X \to \R\) be a real-valued function on the cotangent
bundle that is smooth and homogeneous of degree 0 in the fiber
variables. The H\"ormander space \(H_{(s)}(X)\) consists of
distributions \(u\) on \(X\) such that
\[
\|u\|_{(s)}^2 = \int_{T^*X} (1+|\xi|^2)^{s(x,\xi)} |\hat u(\xi)|^2
\,dx\,d\xi < \infty,
\]
in a suitable local coordinate system. Equivalently,
\(u \in H_{(s)}\) if the pseudodifferential operator with symbol
\((1+|\xi|^2)^{s(x,\xi)/2}\) maps \(u\) into \(L^2\).
\end{definition}

These spaces are natural for describing the regularity of solutions to
operators whose symbol changes type. For example, the Tricomi operator
\(x_1 \partial_1^2 + \partial_2^2\) (which changes from elliptic to
hyperbolic across \(x_1 = 0\)) has regularity properties that are best
described using H\"ormander spaces with a variable order \(s(x,\xi)\)
that depends on the position relative to the sonic line.

The general theory of H\"ormander spaces includes embedding theorems,
trace theorems, and interpolation properties analogous to those for
standard Sobolev spaces, but with the flexibility to capture the
microlocal regularity of solutions near points where the operator
degenerates.

\subsection{The H\"ormander Condition for Hyperbolic Operators}

For hyperbolic operators, H\"ormander established the sharp conditions
for well-posedness of the Cauchy problem. Consider a hyperbolic operator
of order \(m\),
\[
P(t,x, D_t, D_x) = D_t^m + \sum_{j=1}^m a_j(t,x,D_x) D_t^{m-j},
\]
where the \(a_j\) are pseudodifferential operators of order \(j\).

\begin{condition}[Levi Condition]
For the Cauchy problem for a hyperbolic operator with smooth
coefficients to be well-posed in \(C^\infty\), the lower-order terms
must satisfy a condition (the Levi condition) relative to the
characteristic roots. If the characteristic roots \(\lambda_j\) have
multiplicities, the subprincipal symbol must vanish on the
characteristic set to sufficiently high order.
\end{condition}

H\"ormander's analysis revealed the precise relationship between the
multiplicity of the characteristic roots and the number of derivatives
that can be lost in the energy estimates. This work completed the
theory of hyperbolic equations that had been developed by
Petrowsky, Leray, and others.

\subsection{The General Theory of Linear PDEs}

The crowning achievement of H\"ormander's work is the unified treatment
of linear PDEs that emerged from his four-volume treatise. The
organizing principle is the \textit{microlocal} viewpoint: every linear
PDE can be analyzed by examining its symbol in the cotangent bundle,
and the fundamental objects are not the operators themselves but their
wave front relations.

\begin{definition}[Wave Front Relation of an Operator]
For a linear operator \(A: \mathcal{D}'(Y) \to \mathcal{D}'(X)\), the
\textbf{wave front relation} \(\WF'(A) \subset T^*X \times T^*Y\) is
the set of pairs \(((x,\xi), (y,\eta))\) such that \((x,\xi)\) is in
the wave front set of the distribution kernel of \(A\) at \((x,y)\)
with respect to the variable \(\xi\) and \(-\eta\). For a
pseudodifferential operator, \(\WF'(A)\) is contained in the diagonal.
For a Fourier integral operator, \(\WF'(A)\) is the canonical relation.
\end{definition}

The calculus of wave front relations allows one to compute the
microlocal effect of compositions, inverses, and adjoints of linear
operators. This is the language in which the propagation of
singularities theorem is most naturally expressed: if \(P\) is a
pseudodifferential operator and \(A\) is a microlocal parametrix, then
\(\WF'(A)\) is precisely the flow-out of the characteristic set under
the Hamilton vector field.

\section{Mathematical Theory}

\subsection{Microlocal Parametrices}

The fundamental construction in microlocal analysis is the parametric
--- an approximate inverse that is a pseudodifferential or Fourier
integral operator. For an elliptic operator \(P\) of order \(m\), the
parametrix \(Q\) is a pseudodifferential operator of order \(-m\) such
that
\[
Q P = I + R, \qquad P Q = I + S,
\]
where \(R\) and \(S\) are smoothing operators (their distribution
kernels are smooth). The existence of such a parametrix proves elliptic
regularity: if \(Pu \in C^\infty\), then \(u = Q(Pu) - Ru \in C^\infty\),
since \(Q\) preserves smoothness and \(R\) is smoothing.

For operators of real principal type, the parametrix is a Fourier
integral operator whose canonical relation is the bicharacteristic
flow. The construction proceeds by:

\begin{enumerate}
\item Factor the operator into first-order factors near a characteristic
point, using the Malgrange preparation theorem.
\item Solve a transport equation along the bicharacteristics.
\item Sum the resulting oscillatory integrals using the stationary
phase method.
\end{enumerate}

The parametrix construction is the mathematical implementation of the
WKB method: the phase function satisfies the eikonal equation
\(p_m(x, d\phi) = 0\), and the amplitude satisfies the transport
equation \(T a = 0\), where \(T\) is a first-order differential operator
along the bicharacteristic flow.

\subsection{Sharp Regularity: H\"ormander's Theorems}

The sharpest results on the regularity of solutions to linear PDEs are
expressed in terms of Sobolev spaces \(H^s\). For elliptic operators of
order \(m\),
\[
u \in H^s_{\text{loc}} \;\Longrightarrow\; Pu \in H^{s-m}_{\text{loc}},
\]
and conversely, if \(Pu \in H^{s-m}_{\text{loc}}\) and \(u \in
H^{-\infty}_{\text{loc}}\), then \(u \in H^{s}_{\text{loc}}\).

For operators of real principal type, the regularity propagation
theorem takes the form: if \(Pu \in H^{s-m+1}_{\text{loc}}\) and
\(u \in H^t_{\text{loc}}\) for some \(t\), then
\(u \in H^{s}_{\text{loc}}\) microlocally along bicharacteristics, with
a loss of regularity depending on the geometry of the characteristic
set.

\begin{theorem}[Sharp Propagation, H\"ormander]
Let \(P\) be a pseudodifferential operator of order \(m\) with real
principal symbol \(p_m\). Suppose that \(u \in \mathcal{D}'(X)\)
satisfies \(Pu \in H^{s}_{\text{loc}}\). Then
\[
\WF^{s+m-1}(u) \setminus \WF^{s}(Pu) \subset \Sigma,
\]
and this set is invariant under the Hamilton flow,
where \(\WF^{s}(u)\) denotes the \(H^s\)-wave front set (the set of
directions in which \(u\) is not microlocally in \(H^s\)).
\end{theorem}

This theorem gives the exact loss of derivatives: one derivative is
lost relative to the elliptic case, reflecting the fact that
non-elliptic operators cannot fully control the regularity in all
directions simultaneously.

\subsection{The H\"ormander--Weinstein Parametrix}

A particularly elegant construction is the
H\"ormander--Weinstein parametrix for the Cauchy problem for strictly
hyperbolic equations. Consider
\[
\partial_t^2 u - c(x)^2 \Delta u = 0, \quad u(0) = u_0,
\partial_t u(0) = u_1.
\]
The solution operator can be written as a sum of two Fourier integral
operators:
\[
u(t) = A_+(t) u_0 + A_-(t) u_0 + B_+(t) u_1 + B_-(t) u_1,
\]
where each operator has phase function \(\phi_\pm(t,x,\xi) = x\cdot\xi
\pm t c(x)|\xi|\) and amplitude satisfying a transport equation.

The key observation is that the contribution from each characteristic
root can be treated separately, and the total solution is the sum of
these microlocal contributions. This decomposition is the rigorous
version of the method of characteristics for linear hyperbolic
equations and is the starting point for the modern theory of
diffraction and parametrices for non-smooth coefficients.

\subsection{Pseudodifferential Boundary Problems}

H\"ormander also made fundamental contributions to the theory of
boundary value problems for pseudodifferential operators. The
classical theory of elliptic boundary problems (the
Shapiro--Lopatinskii condition) was extended to systems and to
non-elliptic operators.

\begin{definition}[Shapiro--Lopatinskii Condition]
Let \(P\) be an elliptic operator on a domain \(\Omega\) with boundary
\(\partial\Omega\). A boundary operator \(B\) satisfies the
\textbf{Shapiro--Lopatinskii condition} if for every boundary point
and every interior direction, the ordinary differential equation
obtained by freezing coefficients and Fourier-transforming in the
tangential directions has a unique bounded solution satisfying the
transformed boundary condition.
\end{definition}

H\"ormander systematized the theory of pseudodifferential operators on
manifolds with boundary, showing that the boundary value problem can be
reduced to a pseudodifferential equation on the boundary (the
Calder\'on projector), and that the Shapiro--Lopatinskii condition is
equivalent to the ellipticity of this boundary operator.

\subsection{The Method of Stationary Phase}

A crucial technical tool in H\"ormander's work is the method of
stationary phase, which provides the asymptotic expansion of
oscillatory integrals.

\begin{theorem}[Stationary Phase, H\"ormander]
Let \(\phi\) be a smooth, real-valued phase function with a
non-degenerate critical point at \(x_0\) (so \(\nabla\phi(x_0) = 0\)
and \(\det \phi''(x_0) \neq 0\)), and let \(a\) be a smooth amplitude
with compact support. Then as \(\lambda \to \infty\),
\[
\int e^{i\lambda \phi(x)} a(x) \,dx \sim
\left(\frac{2\pi}{\lambda}\right)^{n/2}
\frac{e^{i\lambda \phi(x_0) + i\pi\sigma/4}}{\sqrt{|\det\phi''(x_0)|}}
\sum_{j=0}^\infty \lambda^{-j} L_j a(x_0),
\]
where \(\sigma\) is the signature of \(\phi''(x_0)\) and \(L_j\) are
linear differential operators of order \(2j\).
\end{theorem}

This expansion is the workhorse of microlocal analysis. It is used to:
\begin{itemize}
    \item Compute the composition of pseudodifferential operators.
    \item Derive the asymptotics of Fourier integral operators.
    \item Estimate the kernel of parametrices near the diagonal.
    \item Prove the propagation of singularities via the
    construction of microlocal parametrices.
\end{itemize}

\section{Impact}

\subsection{The Foundation of Microlocal Analysis}

H\"ormander's work created the field of microlocal analysis, which is
now the standard language for the study of linear PDEs. Every
contemporary paper in the analysis of PDEs uses the vocabulary he
introduced: wave front sets, pseudodifferential operators, Fourier
integral operators, symbol classes, and microlocal regularity.

The theory has been extended to:
\begin{itemize}
    \item Complex analysis: the microlocal theory of the
    \(\overline{\partial}\)-operator and CR manifolds.
    \item Spectral theory: the study of eigenvalues and resonances
    via microlocal methods.
    \item Nonlinear PDEs: paradifferential operators (Bony) and
    microlocal analysis of nonlinear equations.
    \item Geometry: the index theorem for elliptic operators and
    the analysis of the Dirac operator.
    \item Inverse problems: boundary rigidity, electrical impedance
    tomography, and the Calder\'on problem.
\end{itemize}

\subsection{The Four-Volume Treatise}

H\"ormander's \textit{The Analysis of Linear Partial Differential
Operators} (Springer, four volumes, 1983--1985) is universally
regarded as the definitive reference. It has been described as
``the Bible of linear PDE theory'' and remains the standard citation
for every major result in the field.

\begin{itemize}
    \item \textbf{Volume 1: Distribution Theory and Fourier Analysis.}
    The foundational volume covers the theory of distributions,
    Fourier transforms, Sobolev spaces, and the basics of linear PDEs.
    It contains the classic exposition of the theory of distributions
    with H\"ormander's signature precision.
    \item \textbf{Volume 2: Differential Operators with Constant
    Coefficients.} The second volume treats constant-coefficient
    operators completely: fundamental solutions, existence theorems,
    the Ehrenpreis--Malgrange theorem on the existence of fundamental
    solutions, and the theory of convolution equations.
    \item \textbf{Volume 3: Pseudodifferential Operators.} The third
    volume develops the calculus of pseudodifferential operators,
    including the symbol classes \(S^m_{\rho,\delta}\), the sharp
    G\r{a}rding inequality, the Fefferman--Phong inequality, and the
    spectral theory of elliptic operators.
    \item \textbf{Volume 4: Fourier Integral Operators.} The final
    volume presents the theory of Fourier integral operators, the
    propagation of singularities, the local solvability condition,
    and the Cauchy problem for hyperbolic equations.
\end{itemize}

The treatise is remarkable for its completeness and for the uniformity
of its presentation: every result is proved from first principles, and
the exposition follows a single, coherent point of view.

\subsection{Applications to Other Fields}

H\"ormander's techniques have found applications far beyond their
original context:

\begin{itemize}
    \item \textbf{Seismic imaging:} Fourier integral operators are used
    in the Kirchhoff migration method for imaging the Earth's interior
    from seismic data.
    \item \textbf{Medical imaging:} The theory of the Radon transform
    and its generalizations uses H\"ormander's results on Fourier
    integral operators and the analysis of singularities.
    \item \textbf{Quantum mechanics:} The semiclassical limit of quantum
    mechanics is described by Fourier integral operators and the
    WKB method, both systematized by H\"ormander.
    \item \textbf{Optics:} Geometrical optics is the leading-order
    approximation of Fourier integral operators, and diffraction
    theory corresponds to the lower-order terms.
    \item \textbf{Control theory:} The H\"ormander condition on vector
    fields appears in the study of controllability and the
    Chow--Rashevskii theorem on accessibility of non-holonomic systems.
\end{itemize}

\subsection{H\"ormander's Influence on Mathematics}

H\"ormander was not only a great researcher but also a great teacher.
His monographs and textbooks shaped generations of analysts:
\begin{itemize}
    \item \textit{Linear Partial Differential Operators} (1963) was the
    first systematic treatment of the subject and won the Fields Medal
    in 1962 (H\"ormander was one of the youngest Fields medalists).
    \item The four-volume treatise (1983--1985) remains the definitive
    reference for the entire field.
    \item His lecture notes and seminar presentations set the standard
    for clarity and rigor.
\end{itemize}

Many of H\"ormander's students and collaborators became leading figures
in analysis, including Richard Melrose, Andr\'as Vasy, and Maciej
Zworski. The H\"ormander tradition in microlocal analysis continues at
Lund, Stanford, MIT, and Berkeley.

\subsection{Undergraduate Education}

While H\"ormander's work is at the research frontier, many of his ideas
now appear in the advanced undergraduate curriculum:
\begin{itemize}
    \item The concept of pseudodifferential operators is introduced in
    courses on Fourier analysis and linear PDEs.
    \item The method of stationary phase is taught in asymptotics and
    applied mathematics courses.
    \item Wave front sets and the propagation of singularities appear
    in advanced PDE courses and mathematical physics.
    \item The H\"ormander condition for hypoellipticity is covered in
    courses on stochastic analysis and control theory.
    \item H\"ormander spaces are used in the theory of function spaces
    and interpolation theory.
\end{itemize}

\section{Frequently Asked Questions}

\subsection*{Q1: What is microlocal analysis?}

Microlocal analysis is the study of distributions and partial
differential equations in the cotangent bundle \(T^*X\) of a manifold
\(X\), where points represent both a position \(x\) and a direction
(frequency) \(\xi\). The key insight is that many properties of PDEs
(ellipticity, hyperbolicity, propagation of singularities) are
\textit{microlocal}: they depend on the behavior of the operator's
symbol in a conic neighborhood of a point in \(T^*X\). The wave front
set \(\WF(u)\) is the microlocal version of the singular support, and
the propagation theorem states that singularities move along
bicharacteristics --- the integral curves of the Hamilton vector field
of the principal symbol.

\subsection*{Q2: What is the H\"ormander condition for local
solvability?}

The H\"ormander condition characterizes when a linear PDE
\(Pu = f\) has a local solution for every smooth right-hand side \(f\).
It requires that the principal symbol \(p_m(x,\xi)\) does not change
sign when the frequency variable crosses the characteristic surface in
a complex direction. Lewy's famous non-solvable operator fails this
condition: its symbol changes sign, which causes the obstruction to
solvability. For operators with real principal symbols (including all
classical types), the condition holds automatically, explaining why
these operators are always locally solvable.

\subsection*{Q3: How are pseudodifferential operators and Fourier
integral operators related?}

A pseudodifferential operator is a special case of a Fourier integral
operator where the phase function is the linear phase
\(\phi(x,\xi) = x \cdot \xi\). Pseudodifferential operators have the
property that their kernels are smooth away from the diagonal; they
correspond to canonical relations that are contained in the diagonal of
\(T^*X \times T^*Y\). Fourier integral operators generalize this by
allowing arbitrary non-degenerate phase functions; their canonical
relations are Lagrangian submanifolds of the cotangent bundle. The
solution operator for the wave equation is a Fourier integral operator
(not a pseudodifferential operator) because it propagates singularities
along non-trivial bicharacteristics.

\subsection*{Q4: What does the propagation of singularities theorem
say intuitively?}

The theorem says that if a function \(u\) satisfies a linear PDE
\(Pu = f\) where \(f\) is smooth, then any singularity of \(u\) must
travel along the characteristic curves of the operator. For the wave
equation, this says that disturbances travel along light rays. For the
Schr\"odinger equation, singularities propagate along classical
trajectories determined by the Hamiltonian \(|\xi|^2\). The theorem is
the rigorous mathematical formulation of the fact that
high-frequency waves follow rays, which is the fundamental principle of
geometrical optics.

\subsection*{Q5: What is the H\"ormander condition for
hypoellipticity?}

An operator \(P\) is hypoelliptic if \(Pu \in C^\infty\) forces
\(u \in C^\infty\). The H\"ormander condition gives a sufficient
condition for operators of the form \(\sum X_j^2 + X_0 + c\), where
\(X_j\) are vector fields: if the Lie algebra generated by the vector
fields has full rank at every point (the brackets of the vector fields
span the tangent space), then the operator is hypoelliptic. This
explains why the heat equation and the Kolmogorov equation are
hypoelliptic even though they are not elliptic: the vector fields
\(\partial_x\) and \(x\partial_y\) have a non-trivial commutator
\(\partial_y\) that supplies the missing direction.

\subsection*{Q6: What are H\"ormander spaces?}

H\"ormander spaces are Sobolev spaces with variable order of
regularity. While the standard Sobolev space \(H^s\) measures
regularity by a fixed number of derivatives, a H\"ormander space
\(H_{(s)}\) allows the number of derivatives to vary with position and
direction. This is natural for operators that change type (like the
Tricomi equation, which is elliptic on one side of a curve and
hyperbolic on the other) or for operators whose symbol degenerates in
specific directions.

\subsection*{Q7: What are the main open problems related to
H\"ormander's work?}

Several deep problems remain:
\begin{enumerate}
    \item \textbf{The strong H\"ormander condition:} Characterizing
    hypoellipticity for general pseudodifferential operators (not just
    sums of squares) in terms of the principal and subprincipal symbols.
    \item \textbf{The Lax--Mizohata condition:} Understanding the
    precise relationship between the H\"ormander condition for local
    solvability and the existence of smooth solutions for non-linear
    equations.
    \item \textbf{Nonlinear microlocal analysis:} Extending the full
    machinery of Fourier integral operators to nonlinear PDEs (the
    program of paradifferential operators and microlocal analysis for
    non-smooth coefficients).
    \item \textbf{Global propagation:} Understanding the propagation
    of singularities on manifolds with boundary and on non-compact
    manifolds, particularly for operators with trapping.
    \item \textbf{Spectral asymptotics:} Sharp estimates on the
    distribution of eigenvalues of non-self-adjoint pseudodifferential
    operators.
\end{enumerate}

\section{Citations}

\subsection*{Primary Sources}
\begin{enumerate}
    \item L.\ H\"ormander, \textit{Linear Partial Differential
    Operators}, Springer, 1963. (Fields Medal monograph.)
    \item L.\ H\"ormander, \textit{Pseudo-differential operators and
    hypoelliptic equations}, Proc.\ Symp.\ Pure Math.\ 10 (1967),
    138--183. (The foundational paper on pseudodifferential operators.)
    \item L.\ H\"ormander, \textit{Fourier integral operators I},
    Acta Math.\ 127 (1971), 79--183. (The creation of Fourier integral
    operator theory.)
    \item L.\ H\"ormander, \textit{The Analysis of Linear Partial
    Differential Operators}, Volumes I--IV, Springer, 1983--1985.
    \item L.\ H\"ormander, \textit{On the existence and the regularity
    of solutions of linear pseudo-differential equations}, Enseign.
    Math.\ 17 (1971), 99--163. (The propagation of singularities.)
    \item L.\ H\"ormander, \textit{Hypoelliptic second order
    differential equations}, Acta Math.\ 119 (1967), 147--171.
    (The sum-of-squares theorem.)
    \item L.\ H\"ormander, \textit{Differential equations without
    solutions}, Math.\ Ann.\ 140 (1960), 169--173. (The local
    solvability condition.)
\end{enumerate}

\subsection*{Biographies and Surveys}
\begin{enumerate}
    \item S.\ M.\ Salamon, \textit{Lars H\"ormander (1931--2012)},
    Biographical Memoirs of Fellows of the Royal Society 60 (2014),
    235--256.
    \item R.\ B.\ Melrose, \textit{The work of Lars H\"ormander},
    Notices AMS 60 (2013), 720--725.
    \item M.\ Zworski, \textit{Lars H\"ormander --- a biography},
    European Math.\ Soc.\ Newsletter 89 (2013), 18--22.
    \item R.\ B.\ Melrose, \textit{Microlocal analysis and the work of
    Lars H\"ormander}, Proc.\ Int.\ Cong.\ Math.\ 2014, 69--86.
\end{enumerate}

\subsection*{Textbooks and Expository Works}
\begin{enumerate}
    \item M.\ E.\ Taylor, \textit{Pseudodifferential Operators},
    Princeton University Press, 1981.
    \item F.\ Treves, \textit{Introduction to Pseudodifferential and
    Fourier Integral Operators}, Plenum Press, 1980.
    \item J.\ J.\ Duistermaat, \textit{Fourier Integral Operators},
    Birkh\"auser, 1996.
    \item A.\ Grigis and J.\ Sj\"ostrand, \textit{Microlocal Analysis
    for Differential Operators}, Cambridge University Press, 1994.
    \item M.\ Zworski, \textit{Semiclassical Analysis}, AMS, 2012.
    \item L.\ C.\ Evans, \textit{Partial Differential Equations},
    2nd ed., AMS, 2010.
    \item G.\ B.\ Folland, \textit{Introduction to Partial Differential
    Equations}, 2nd ed., Princeton University Press, 1995.
    \item S.\ Alinhac and P.\ G\'erard, \textit{Pseudo-differential
    Operators and the Nash--Moser Theorem}, AMS, 2007.
    \item N.\ Lerner, \textit{Metrics on the Phase Space and
    Non-Selfadjoint Pseudo-Differential Operators}, Birkh\"auser, 2010.
\end{enumerate}

\section{Timeline}

\begin{tabular}{|l|l|}
\hline
\textbf{Year} & \textbf{Event} \\ \hline
1931 & Born on January 24 in Mj\"allby, Blekinge, Sweden \\ \hline
1955 & Ph.D.\ in Mathematics, Lund University (advisor: Marcel Riesz) \\ \hline
1956--1957 & Postdoctoral visits to Stockholm and the United States \\
\hline
1957 & Appointed Docent at Lund University \\ \hline
1960 & H\"ormander condition for local solvability published \\
\hline
1962 & Awarded the Fields Medal at the ICM in Stockholm \\
& (the youngest Fields medalist at age 31) \\ \hline
1963 & Published \textit{Linear Partial Differential Operators} \\
\hline
1963--1965 & Visiting Professor at Stanford University and the Institute
for Advanced Study, Princeton \\ \hline
1964 & Appointed Professor at the University of Lund \\ \hline
1967 & H\"ormander's hypoellipticity theorem (sums of squares) \\
\hline
1968 & Elected to the Royal Swedish Academy of Sciences \\ \hline
1970--1971 & Theory of Fourier integral operators and propagation
of singularities \\ \hline
1971 & Pseudodifferential operator theory fully developed \\ \hline
1974 & Elected to the French Academy of Sciences \\ \hline
1975 & H\"ormander's results on the Levi condition for hyperbolic
operators \\ \hline
1979 & Elected to the American Academy of Arts and Sciences
(foreign honorary member) \\ \hline
1980s & Continued work on the spectral theory of pseudodifferential
operators \\ \hline
1983--1985 & Publication of \textit{The Analysis of Linear Partial
Differential Operators} (four volumes) \\
\hline
1988 & Awarded the Wolf Prize in Mathematics (shared with
Friedrich Hirzebruch) \\ \hline
1996 & Retired from the University of Lund; appointed Professor
Emeritus \\ \hline
2006 & Awarded the Steele Prize for Mathematical Exposition
(American Mathematical Society) \\ \hline
2012 & Died on November 25 in Lund, Sweden, at age 81 \\ \hline
\end{tabular}

\section*{Exercises}
\addcontentsline{toc}{section}{Exercises}

\begin{enumerate}
    \item Let \(u\) be the characteristic function of the unit ball in
    \(\R^3\). Compute \(\WF(u)\). Show that the singularities occur
    only at the boundary and that the wave front set consists of
    conormal directions to the sphere.
    \item Prove that for the Dirac delta \(\delta_0\) on \(\R^n\),
    \(\WF(\delta_0) = \{(0,\xi): \xi \neq 0\}\). Show that if a
    distribution has \(\WF(u) = \varnothing\), then \(u \in C^\infty\).
    \item Let \(P = -\partial_x^2 - \partial_y^2\) be the Laplacian on
    \(\R^2\). Compute the principal symbol \(p_2(x,y,\xi,\eta)\) and
    the Hamilton vector field \(H_{p_2}\). Show that the
    bicharacteristics are the straight lines \(\xi =\) constant,
    \(\eta =\) constant.
    \item For the wave operator \(\square = \partial_t^2 - \partial_x^2\)
    on \(\R^2\), compute the principal symbol and the bicharacteristics.
    Show that the propagation of singularities theorem recovers the
    fact that solutions travel along light cones.
    \item Verify that Lewy's operator
    \[
    L = \frac{\partial}{\partial x_1} + i\frac{\partial}{\partial x_2}
    - 2i(x_1 + i x_2) \frac{\partial}{\partial x_3}
    \]
    has principal symbol \(p_1(x,\xi) = \xi_1 + i\xi_2 - 2i(x_1 + i x_2)\xi_3\).
    Show that the H\"ormander sign condition fails at the origin,
    explaining why \(L\) is not locally solvable.
    \item Let \(X_1 = \partial_x\) and \(X_2 = x\partial_y\) on \(\R^2\).
    Compute their commutator \([X_1, X_2]\). Show that the operator
    \(P = X_1^2 + X_2^2\) satisfies H\"ormander's condition and is
    therefore hypoelliptic.
    \item Prove that if \(a \in S^{m_1}\) and \(b \in S^{m_2}\), then
    the composition symbol \(c(x,\xi) = a(x,D) b(x,\xi)\) (where
    \(b(x,\xi)\) acts on \(a\)) belongs to \(S^{m_1+m_2}\) and has
    leading term \(a(x,\xi) b(x,\xi)\).
    \item Implement the stationary phase method for the one-dimensional
    integral
    \[
    I(\lambda) = \int_{-1}^1 e^{i\lambda x^2} \,dx.
    \]
    Compute the asymptotic expansion as \(\lambda \to \infty\) and
    compare with the exact evaluation in terms of Fresnel integrals.
    \item Prove that the heat operator \(\partial_t - \partial_x^2\)
    on \(\R^2\) is hypoelliptic by verifying H\"ormander's condition
    for the vector fields \(X_0 = \partial_t\) and \(X_1 = \partial_x\).
    \item Let \(u\) be a distribution and let \(\phi \in C_0^\infty\).
    Show that \(\WF(\phi u) \subset \WF(u)\) and that
    \(\WF(u) = \bigcup_\phi \WF(\phi u)\) (where the union is over a
    partition of unity).
    \item For the Schr\"odinger operator \(i\partial_t + \Delta\) on
    \(\R^{1+n}\), compute the principal symbol and the bicharacteristics.
    Show that singularities propagate along straight lines in the
    spatial variables.
    \item Derive the composition formula for two pseudodifferential
    operators whose symbols are independent of \(x\) (so-called
    Fourier multipliers). Show that in this case the composition is
    exact (the asymptotic expansion terminates at the first term).
    \item Let \(P\) be an elliptic operator on a compact manifold.
    Construct a parametrix \(Q\) of order \(-m\) such that
    \(P Q = I + R\) where \(R\) has a smooth kernel. Use this to prove
    the elliptic regularity theorem.
    \item Show that H\"ormander's condition for local solvability is
    equivalent to the condition that the principal symbol \(p_m\)
    does not vanish on the characteristic set to infinite order when
    the frequency is perturbed in a complex direction normal to the
    characteristic surface.
    \item (Project) Read H\"ormander's original 1971 paper on Fourier
    integral operators and write a summary of the main theorems. Focus
    on the construction of the parametrix for operators of real
    principal type.
\end{enumerate}

\section*{Glossary}
\addcontentsline{toc}{section}{Glossary}

\begin{description}
    \item[Bicharacteristic] An integral curve of the Hamilton vector
    field \(H_{p_m}\) on the characteristic set \(\Sigma = \{p_m = 0\}\).
    \item[Canonical relation] A Lagrangian submanifold of
    \(T^*X \times T^*Y\) describing the wave front set of the kernel of
    a Fourier integral operator.
    \item[Characteristic set] The set \(\{(x,\xi) \in T^*X \setminus 0
    : p_m(x,\xi) = 0\}\) where the principal symbol vanishes.
    \item[Elliptic operator] An operator whose principal symbol vanishes
    only at \(\xi = 0\); equivalently, \(|p_m(x,\xi)| \geq C|\xi|^m\)
    for large \(|\xi|\).
    \item[Fourier integral operator] An operator defined by an
    oscillatory integral with a non-degenerate phase function,
    generalizing pseudodifferential operators.
    \item[Hamilton vector field] The vector field
    \(H_p = \sum (\partial_{\xi_j} p \,\partial_{x_j} - \partial_{x_j} p
    \,\partial_{\xi_j})\).
    \item[Hypoelliptic operator] An operator \(P\) such that
    \(Pu \in C^\infty\) implies \(u \in C^\infty\).
    \item[H\"ormander condition] The condition on iterated commutators
    of vector fields that guarantees hypoellipticity of
    \(\sum X_j^2 + X_0\).
    \item[H\"ormander space] A Sobolev space with variable order of
    regularity, denoted \(H_{(s)}\) where \(s = s(x,\xi)\).
    \item[Levi condition] A condition on the lower-order terms of a
    hyperbolic operator ensuring well-posedness of the Cauchy problem.
    \item[Local solvability] The property that \(Pu = f\) has a
    solution for every smooth \(f\) in a neighborhood.
    \item[Microlocal analysis] The study of distributions and PDEs in
    the cotangent bundle, analyzing singularities by position and
    frequency simultaneously.
    \item[Parametrix] An approximate inverse of a differential or
    pseudodifferential operator, modulo smoothing operators.
    \item[Principal symbol] The highest-order part of the symbol of a
    differential operator, capturing its essential microlocal behavior.
    \item[Propagation of singularities] The theorem that \(\WF(u)\)
    is invariant under the Hamilton flow for solutions of \(Pu = f\)
    with \(f\) smooth.
    \item[Pseudodifferential operator] An operator defined by
    \((2\pi)^{-n} \int e^{ix\cdot\xi} a(x,\xi) \hat u(\xi) \,d\xi\)
    where the symbol \(a(x,\xi)\) is a smooth function of moderate
    growth.
    \item[Shapiro--Lopatinskii condition] The necessary and sufficient
    condition for an elliptic boundary value problem to be Fredholm.
    \item[Stationary phase] The asymptotic expansion of oscillatory
    integrals \(\int e^{i\lambda\phi(x)} a(x)\,dx\) as
    \(\lambda \to \infty\) in terms of critical points of \(\phi\).
    \item[Sum of squares] An operator of the form
    \(\sum_{j=1}^k X_j^2 + X_0 + c\), studied by H\"ormander for
    hypoellipticity.
    \item[Symbol class] The space \(S^m_{\rho,\delta}\) of symbols
    satisfying symbol estimates with decay order \(m\) and smoothness
    indices \(\rho, \delta\).
    \item[Wave front set] The set \(\WF(u) \subset T^*X \setminus 0\)
    recording the singular directions of a distribution \(u\) at each
    point.
    \item[WKB method] The Wentzel--Kramers--Brillouin approximation,
    a semiclassical method for constructing approximate solutions of
    PDEs using oscillatory exponentials and transport equations.
\end{description}
